\documentclass[../main.tex]{subfiles}
\begin{document}

\section{Lecture 22 -- Sound Waves in Fluids}\label{sec:lec22}
\begin{center}
  \textbf{Date:} June 22, 2026
\end{center}

\subsection*{Overview}
This lecture closes the discussion of Stokes flow past a sphere with two useful
clarifications, then turns to a new phenomenon: \textbf{small disturbances in a
fluid at rest}. The physical example is familiar from everyday life. A hammer
tap, a lightning strike, a vibrating surface, or a droplet falling into water
creates a local disturbance; the disturbance propagates away as a wave. In air
we call it sound.

The key idea is that a sound wave in a fluid is a coupled oscillation of
\textbf{pressure}, \textbf{density}, and \textbf{velocity}. Unlike an elastic
solid, a fluid at rest has no shear modulus, so it cannot support transverse
sound waves. The only propagating acoustic mode is longitudinal: the fluid
velocity oscillates parallel to the direction in which the wave travels.

The mathematical strategy is the same perturbation method used throughout the
course. Start from a homogeneous static solution
\[
  \vel=\bvec0,\qquad p=p_0,\qquad \rho=\rho_0,
\]
add small perturbations, keep only first order terms, and combine the linearized
Euler equation, the linearized continuity equation, and the equation of state.
The result is a wave equation with sound speed
\[
  c^2=\left(\pdv{p}{\rho}\right)_s,
\]
where the derivative is adiabatic/isentropic because sound oscillations are fast
enough that heat has no time to flow appreciably between neighboring compressed
and rarefied regions.

\medskip\noindent\textbf{Topics:}
\begin{enumerate}
  \item Clarifications on the Stokes-sphere solution: multipoles, radial stream
    function, and why the biharmonic source term vanishes.
  \item Flow regimes around a sphere as the Reynolds number increases.
  \item Phenomenology and historical measurements of sound in air and water.
  \item Comparison with elastic waves in solids.
  \item Linearization around a fluid at rest.
  \item Absence of transverse sound waves in an ideal fluid.
  \item Longitudinal velocity potential and the acoustic wave equation.
  \item Sound speed in an ideal gas and in water.
\end{enumerate}

%%─────────────────────────────────────────────────────────────────────────────
\subsection{Two Clarifications on Stokes Flow Past a Sphere}\label{sec:lec22:stokes-clarifications}
%%─────────────────────────────────────────────────────────────────────────────

\subsubsection*{Conceptual Background}
The previous lecture derived Stokes' law for a sphere moving slowly through a
viscous fluid. The result was
\begin{equation}
  \boxed{\;\bvec F_D=6\pi\eta R\,\bvec V_0\;}
  \label{eq:lec22_stokes_drag_recall}
\end{equation}
for a sphere of radius $R$ in a fluid of dynamic viscosity $\eta$, with far-field
velocity $\bvec V_0$ in the sphere frame. The corresponding flow is called
\textbf{creeping flow} (or Stokes flow): it is the low-Reynolds-number regime
\[
  \reynolds=\frac{V_0R}{\nu}\ll1,\qquad \nu=\eta/\rho.
\]
The professor returned to two points in the derivation that deserve to be made
explicit.

\subsubsection*{Why the scalar function may be radial}
In \cref{sec:lec21} the velocity disturbance was written using one scalar
function $f(r)$,
\begin{equation}
  \vel=\bvec V_0+\grad\times\grad\times\big(f(r)\bvec V_0\big).
  \label{eq:lec22_stokes_stream_recall}
\end{equation}
The question is: why is $f$ allowed to depend only on $r$, and not also on
$\theta$?

The answer is not that the whole physical setup is spherically symmetric. It is
not: the vector $\bvec V_0$ selects a preferred direction. The actual reason is
\textbf{linearity}. In creeping flow the inertial term
$(\vel\cdot\grad)\vel$ is neglected, so the Stokes equations are linear. Linear
rotationally invariant equations allow a multipole decomposition: each angular
momentum sector $\ell$ obeys its own equation and does not mix with the others,
except through boundary conditions.

The boundary condition at infinity is a uniform vector field,
\[
  \vel\to\bvec V_0\qquad(r\to\infty),
\]
and a vector transforms as the $\ell=1$ representation of rotations. Therefore
the whole solution lives in the $\ell=1$ sector. The angular dependence is
already supplied by the vector $\bvec V_0$ and by the differential operators in
\eqref{eq:lec22_stokes_stream_recall}. Adding an arbitrary explicit
$\theta$-dependence to $f$ would introduce other multipoles, which the boundary
data do not excite.

\nt{The phrase ``$f=f(r)$'' should therefore be read as: after the required
$\ell=1$ angular structure has been factored out through $\bvec V_0$, the
remaining unknown amplitude is radial.}

\subsubsection*{Dipole plus Stokeslet}
The velocity field found in \cref{sec:lec21} can be read as a sum of two
$\ell=1$ vector fields:
\begin{equation}
  \vel=\bvec V_0-\frac{3R}{4}\,
  \frac{\bvec V_0+(\bvec V_0\cdot\uvec r)\uvec r}{r}
  -\frac{R^3}{4}\,
  \frac{\bvec V_0-3(\bvec V_0\cdot\uvec r)\uvec r}{r^3}.
  \label{eq:lec22_stokes_multipoles}
\end{equation}
The last term is the familiar \textbf{potential dipole} field: it can be written
as the gradient of a scalar dipole potential. The $1/r$ term is also an
$\ell=1$ vector field, but it is not the gradient of a scalar potential. It is
the \textbf{Stokeslet} part of the solution, the long-ranged viscous disturbance
associated with the net force on the sphere.

\textit{Significance:} Both pieces carry the same angular momentum label
$\ell=1$, but they are physically different: the $1/r^3$ part is a potential
dipole correction, while the $1/r$ Stokeslet is the viscous force-carrying tail.

\subsubsection*{Why the biharmonic equation is homogeneous}
Another small gap in the derivation was the step from
\[
  \grad(\nabla^4 f)\times\bvec V_0=\bvec0
\]
to $\nabla^4 f=0$. The cross product vanishes if
$\grad(\nabla^4 f)$ is parallel to $\bvec V_0$. But $f=f(r)$ is radial, so
$\grad(\nabla^4 f)$ points in the radial direction $\uvec r$. A radial vector
cannot be parallel to the same fixed vector $\bvec V_0$ at every point in space,
unless it is zero. Hence
\begin{equation}
  \boxed{\;\nabla^4 f=0\;}
  \label{eq:lec22_biharmonic_homogeneous}
\end{equation}
up to an irrelevant constant fixed by the far-field condition.

\textit{Significance:} The homogeneity of the biharmonic equation follows from
the combination of radial dependence and a fixed far-field direction; it is not
an extra physical assumption.

%%─────────────────────────────────────────────────────────────────────────────
\subsection{Flow Regimes Around a Sphere}\label{sec:lec22:reynolds-regimes}
%%─────────────────────────────────────────────────────────────────────────────

\subsubsection*{Physical Motivation}
The exact Stokes solution is derived in the limit $\reynolds\to0$, but the
qualitative picture remains useful for small finite Reynolds number. As
$\reynolds$ grows, the flow around a sphere does not merely change its numerical
coefficients; it changes its \emph{type}.

\begin{itemize}
  \item \textbf{$\reynolds\lesssim1$--$10$: creeping/laminar flow.} The Stokes
  picture is qualitatively correct. The flow is steady, smooth, and has no wake
  vortices.

  \item \textbf{$10\lesssim\reynolds\lesssim150$: steady separated wake.} The
  fore-aft symmetry breaks and a recirculating region appears behind the sphere,
  but the flow can still be steady in time.

  \item \textbf{$150\lesssim\reynolds\lesssim10^3$--$10^4$: vortex shedding.}
  Vortices are shed alternately from the two sides of the body, forming the
  von Karman vortex street. The flow is now time-dependent but still visibly
  organized.

  \item \textbf{Large $\reynolds$: turbulence.} The wake becomes irregular and
  chaotic. One can often describe statistical properties, but not a single
  simple steady streamline pattern.
\end{itemize}

\nt{The precise transition values depend on geometry, roughness, perturbations,
and on whether the body is a sphere or a cylinder. The point for this course is
the qualitative sequence: Stokes flow, steady wake, periodic vortex shedding,
then turbulence.}

\textit{Significance:} High viscosity suppresses complexity and produces
ordered flow. Decreasing viscosity, or increasing speed/size, increases
$\reynolds$ and allows new nonlinear flow structures to appear.

%%─────────────────────────────────────────────────────────────────────────────
\subsection{Sound as a Propagating Fluid Disturbance}\label{sec:lec22:phenomenology}
%%─────────────────────────────────────────────────────────────────────────────

\subsubsection*{Physical Motivation}
Now consider a fluid initially at rest. Strike it, push it, vibrate a boundary,
or create a small pressure pulse. The disturbance propagates through the fluid.
In air, this is ordinary sound. In water, the same mechanism gives underwater
acoustic waves.

The finite speed of sound is visible without special equipment. We see lightning
almost immediately, but hear thunder seconds later. If the delay is about
$10\,\mathrm{s}$, the storm is roughly $3$--$3.5\,\mathrm{km}$ away, because the
speed of sound in air is about
\[
  c_{\mathrm{air}}\simeq 343\,\mathrm{m/s}
\]
near room temperature.

Historically, this speed was measured long before electronic timing. Mersenne
and others used visible flashes and loud sounds such as cannon fire to compare
the arrival time of light and sound. Later measurements improved the value to
the modern range. For water, Colladon and Sturm measured sound propagation
across Lake Geneva in 1826 using an underwater bell and an underwater listening
horn, finding a value close to
\[
  c_{\mathrm{water}}\simeq 1.48\,\mathrm{km/s}.
\]
Thus sound travels roughly five times faster in water than in air.

%%─────────────────────────────────────────────────────────────────────────────
\subsection{Comparison with Elastic Waves in Solids}\label{sec:lec22:solid-comparison}
%%─────────────────────────────────────────────────────────────────────────────

\subsubsection*{Conceptual Background}
In an elastic solid, the dynamical variable is the displacement field
$\bvec u(\curcoord,t)$. Small-amplitude elastic waves split into:
\begin{itemize}
  \item \textbf{longitudinal waves}, where $\bvec u$ is parallel to the
  propagation direction;
  \item \textbf{transverse waves}, where $\bvec u$ is perpendicular to the
  propagation direction. There are two independent transverse polarizations.
\end{itemize}
For an isotropic solid,
\begin{equation}
  \boxed{\;c_P^2=\frac{\lambda+2\mu}{\rho},\qquad
  c_S^2=\frac{\mu}{\rho}\;}
  \label{eq:lec22_elastic_wave_speeds}
\end{equation}
where $\lambda,\mu$ are the Lame constants and $\mu$ is the shear modulus.

\textit{Significance:} The transverse waves exist only because a solid resists
shear. A fluid at rest has no static shear modulus, so the transverse branch has
no fluid analogue.

%%─────────────────────────────────────────────────────────────────────────────
\subsection{Linearization Around a Fluid at Rest}\label{sec:lec22:linearization}
%%─────────────────────────────────────────────────────────────────────────────

\subsubsection*{Mathematical Setup}
Start from a homogeneous static fluid:
\begin{equation}
  \vel_0=\bvec0,\qquad p=p_0,\qquad \rho=\rho_0,
  \label{eq:lec22_background}
\end{equation}
with no important body force. Add small perturbations:
\begin{equation}
  \vel=\delta\vel,\qquad p=p_0+p',\qquad \rho=\rho_0+\rho',
  \label{eq:lec22_perturbations}
\end{equation}
where $|\delta\vel|$, $|p'|$, and $|\rho'|$ are first-order quantities. Below we
write $\vel$ for the perturbation velocity, since the background velocity is
zero.

For acoustic propagation we first neglect viscosity. Viscosity can attenuate a
sound wave, but it is not what creates the wave. The relevant equations are
Euler's equation and mass continuity:
\begin{align}
  \pdv{\vel}{t}+(\vel\cdot\grad)\vel&=-\frac{1}{\rho}\grad p,
  \label{eq:lec22_euler_full}\\
  \pdv{\rho}{t}+\divg(\rho\vel)&=0.
  \label{eq:lec22_continuity_full}
\end{align}

\subsubsection*{Linearized Euler equation}
In Euler's equation, $\partial_t\vel$ is first order, while
$(\vel\cdot\grad)\vel$ is second order and is dropped. Also
\[
  \grad p=\grad(p_0+p')=\grad p',
\]
because $p_0$ is homogeneous. Since $\grad p'$ is already first order, replacing
$1/\rho$ by $1/\rho_0$ is consistent to first order. Thus
\begin{equation}
  \boxed{\;\pdv{\vel}{t}=-\frac{1}{\rho_0}\grad p'\;}
  \label{eq:lec22_linear_euler}
\end{equation}
\textit{Significance:} Pressure gradients accelerate the fluid. This is the
``inertia'' half of sound: a compressed region pushes neighboring fluid outward.

\subsubsection*{Linearized continuity equation}
Expanding continuity,
\[
  \pdv{}{t}(\rho_0+\rho')+\divg\big((\rho_0+\rho')\vel\big)=0.
\]
The background density is constant, and $\rho'\vel$ is second order, so
\begin{equation}
  \boxed{\;\pdv{\rho'}{t}+\rho_0\,\divg\vel=0\;}
  \label{eq:lec22_linear_continuity}
\end{equation}
\textit{Significance:} Converging velocity increases density; diverging velocity
decreases density. This is the ``compression'' half of sound.

\subsubsection*{Linearized equation of state}
The fluid also needs a thermodynamic closure. For small disturbances,
\begin{equation}
  p'= \left(\pdv{p}{\rho}\right)_s \rho'.
  \label{eq:lec22_eos_linear}
\end{equation}
Define
\begin{equation}
  \boxed{\;c^2\equiv\left(\pdv{p}{\rho}\right)_s\;}
  \label{eq:lec22_sound_speed_def}
\end{equation}
where the derivative is taken at constant entropy $s$.

\textit{Significance:} The derivative $c^2$ measures how strongly pressure
responds to compression. It is the restoring force per unit density, exactly the
structure expected for a wave speed.

%%─────────────────────────────────────────────────────────────────────────────
\subsection{Why There Are No Transverse Sound Waves in a Fluid}\label{sec:lec22:no-transverse}
%%─────────────────────────────────────────────────────────────────────────────

\subsubsection*{Physical Motivation}
In a transverse wave, the velocity field has a rotational part: the fluid would
oscillate sideways relative to the propagation direction. A static fluid cannot
support shear stress, so we expect no propagating transverse sound mode.

\subsubsection*{Mathematical Check}
Take the curl of the linearized Euler equation
\eqref{eq:lec22_linear_euler}:
\begin{equation}
  \pdv{}{t}(\curl\vel)
  =-\frac{1}{\rho_0}\curl(\grad p')=\bvec0.
  \label{eq:lec22_vorticity_linear}
\end{equation}
With $\bvec\omega=\curl\vel$,
\begin{equation}
  \boxed{\;\pdv{\bvec\omega}{t}=0\;}
  \label{eq:lec22_no_transverse}
\end{equation}
This is not a wave equation: there is no second time derivative balanced by a
spatial Laplacian. If a rotational perturbation is present in an ideal fluid, it
is simply frozen at this order; with viscosity included, it diffuses and
decays.

\textit{Significance:} Fluids have no acoustic shear branch. Ordinary sound in
air or water is longitudinal.

%%─────────────────────────────────────────────────────────────────────────────
\subsection{Longitudinal Waves and the Acoustic Wave Equation}\label{sec:lec22:wave-equation}
%%─────────────────────────────────────────────────────────────────────────────

\subsubsection*{Mathematical Setup}
The propagating acoustic part is irrotational, so write the velocity perturbation
using a scalar potential:
\begin{equation}
  \vel=\grad\phi.
  \label{eq:lec22_velocity_potential}
\end{equation}
Substitute this into \eqref{eq:lec22_linear_euler}:
\[
  \grad\!\left(\pdv{\phi}{t}\right)=-\frac{1}{\rho_0}\grad p'
  \quad\Longrightarrow\quad
  \grad\!\left(\pdv{\phi}{t}+\frac{p'}{\rho_0}\right)=\bvec0.
\]
The expression in parentheses may depend only on time; this time-dependent
constant can be absorbed into the freedom $\phi\mapsto\phi+g(t)$. Choose the
gauge in which it vanishes:
\begin{equation}
  \boxed{\;p'=-\rho_0\,\pdv{\phi}{t}\;}
  \label{eq:lec22_pressure_potential}
\end{equation}

Now use continuity. Since $\divg\vel=\lapl\phi$,
\begin{equation}
  \pdv{\rho'}{t}+\rho_0\lapl\phi=0.
  \label{eq:lec22_cont_phi}
\end{equation}
The equation of state gives $\rho'=p'/c^2$. With
\eqref{eq:lec22_pressure_potential},
\[
  \rho'=-\frac{\rho_0}{c^2}\pdv{\phi}{t}.
\]
Substituting this into \eqref{eq:lec22_cont_phi},
\[
  -\frac{\rho_0}{c^2}\pdv[2]{\phi}{t}+\rho_0\lapl\phi=0.
\]
Cancel $\rho_0$:
\begin{equation}
  \boxed{\;\lapl\phi=\frac{1}{c^2}\pdv[2]{\phi}{t}\;}
  \label{eq:lec22_wave_phi}
\end{equation}

\textit{Significance:} The potential $\phi$ obeys the ordinary wave equation.
Once $\phi$ is known, the velocity, pressure, and density perturbations are
obtained from
\[
  \vel=\grad\phi,\qquad p'=-\rho_0\phi_t,\qquad \rho'=\frac{p'}{c^2}.
\]

\thm[acoustic-wave]{Acoustic Wave Equation in a Fluid}{%
  Small irrotational perturbations of a homogeneous inviscid fluid at rest obey
  \begin{equation}
    \boxed{\;\lapl\phi-\frac{1}{c^2}\pdv[2]{\phi}{t}=0,\qquad
    c^2=\left(\pdv{p}{\rho}\right)_s\;}
    \label{eq:lec22_acoustic_wave_thm}
  \end{equation}
  with $\vel=\grad\phi$, $p'=-\rho_0\phi_t$, and $\rho'=p'/c^2$.
}
\textit{Significance:} Sound is a pressure-density restoring force coupled to
fluid inertia. The wave speed is fixed by thermodynamics and density, not by
viscosity.

\subsubsection*{Plane-wave solution}
For a monochromatic plane wave traveling in the $z$ direction, take
\begin{equation}
  \phi=\Phi_0\cos(kz-\omega t),\qquad \omega=ck.
  \label{eq:lec22_plane_phi}
\end{equation}
Then
\begin{align}
  v_z&=\pdv{\phi}{z}=-k\Phi_0\sin(kz-\omega t),
  \label{eq:lec22_plane_v}\\
  p'&=-\rho_0\pdv{\phi}{t}=-\rho_0\omega\Phi_0\sin(kz-\omega t),
  \label{eq:lec22_plane_p}\\
  \rho'&=\frac{p'}{c^2}.
  \label{eq:lec22_plane_rho}
\end{align}
Thus, for a wave traveling in the positive $z$ direction,
\begin{equation}
  \boxed{\;p'=\rho_0 c\,v_z,\qquad \rho'=\frac{\rho_0}{c}\,v_z\;}
  \label{eq:lec22_plane_relations}
\end{equation}
with signs understood relative to the propagation direction.

\textit{Significance:} A sound wave consists of alternating compressions and
rarefactions: where $p'$ and $\rho'$ are positive, the fluid is compressed; where
they are negative, it is rarefied. The velocity oscillates longitudinally.

\begin{figure}[h]
\centering
\begin{tikzpicture}[scale=1.0]
  \draw[->,thick] (-0.4,0) -- (8.2,0) node[right] {$z$};
  \draw[myblue,thick,domain=0:7.6,samples=120]
    plot (\x,{0.55*sin(180*\x/1.9)});
  \node[myblue] at (1.5,0.9) {$p',\,\rho'$};
  \foreach \x in {0.45,1.2,1.95,2.7,3.45,4.2,4.95,5.7,6.45,7.2}{
    \draw[vvec] (\x,-1.05) -- ++({0.35*cos(180*\x/1.9)},0);
  }
  \node[vcol] at (1.5,-1.45) {$\vel\parallel\hat{\bvec z}$};
  \node[font=\small] at (2.85,0.35) {compression};
  \node[font=\small] at (4.75,-0.35) {rarefaction};
  \draw[width] (0,-1.85) -- (3.8,-1.85) node[midway,below] {$\lambda$};
\end{tikzpicture}
\caption{A one-dimensional sound wave: pressure and density oscillate between
compression and rarefaction, while the fluid velocity oscillates parallel to the
propagation direction.}
\label{fig:lec22_sound_wave}
\end{figure}

%%─────────────────────────────────────────────────────────────────────────────
\subsection{Sound Speed in Air and Water}\label{sec:lec22:sound-speed}
%%─────────────────────────────────────────────────────────────────────────────

\subsubsection*{Ideal gas}
For an ideal gas undergoing an adiabatic small oscillation,
\[
  p\rho^{-\gamma}=\mathrm{const},
\]
so
\[
  p=K\rho^\gamma,\qquad
  \left(\pdv{p}{\rho}\right)_s=\gamma K\rho^{\gamma-1}
  =\gamma\frac{p}{\rho}.
\]
Therefore
\begin{equation}
  \boxed{\;c_{\mathrm{gas}}=\sqrt{\gamma\frac{p_0}{\rho_0}}\;}
  \label{eq:lec22_ideal_gas_sound}
\end{equation}
For air near room temperature, $\gamma\simeq7/5$ because air is mostly diatomic
nitrogen and oxygen. With $p_0\simeq1.0\times10^5\,\mathrm{Pa}$ and
$\rho_0\simeq1.2\,\mathrm{kg/m^3}$,
\begin{equation}
  c_{\mathrm{air}}\simeq 340\,\mathrm{m/s}.
  \label{eq:lec22_air_number}
\end{equation}

\nt{Newton effectively used the isothermal estimate
$c^2=p_0/\rho_0$, which is too small by a factor $\sqrt{\gamma}$. Laplace's
adiabatic correction supplies the missing factor. For a monatomic gas
$\gamma=5/3$; for ordinary air, $\gamma\simeq7/5$.}

\subsubsection*{Water}
For a liquid it is common to use the adiabatic bulk modulus
\begin{equation}
  B_s\equiv \rho_0\left(\pdv{p}{\rho}\right)_s.
  \label{eq:lec22_bulk_modulus}
\end{equation}
Then
\begin{equation}
  \boxed{\;c_{\mathrm{liquid}}=\sqrt{\frac{B_s}{\rho_0}}\;}
  \label{eq:lec22_liquid_sound}
\end{equation}
For water, $B_s\simeq2.2\,\mathrm{GPa}$ and
$\rho_0\simeq10^3\,\mathrm{kg/m^3}$, so
\begin{equation}
  c_{\mathrm{water}}\simeq
  \sqrt{\frac{2.2\times10^9}{10^3}}\;\mathrm{m/s}
  \simeq1.5\times10^3\,\mathrm{m/s}.
  \label{eq:lec22_water_number}
\end{equation}

\textit{Significance:} Water is much harder to compress than air. Although it is
also denser, the large bulk modulus wins, so sound travels much faster in water.

%%─────────────────────────────────────────────────────────────────────────────
\subsection*{Summary}
%%─────────────────────────────────────────────────────────────────────────────
\begin{itemize}
  \item In Stokes flow past a sphere, $f=f(r)$ because linearity lets us
  decompose into multipoles and the far-field vector $\bvec V_0$ excites only
  the $\ell=1$ sector.
  \item The Stokes velocity contains an $\ell=1$ potential dipole part
  ($1/r^3$) and an $\ell=1$ non-potential Stokeslet part ($1/r$).
  \item Flow past a sphere changes qualitatively with Reynolds number:
  creeping flow $\to$ steady separated wake $\to$ von Karman vortex shedding
  $\to$ turbulence.
  \item A sound wave in a fluid is a small longitudinal disturbance of pressure,
  density, and velocity around a homogeneous rest state.
  \item Linearized Euler:
  $\partial_t\vel=-(1/\rho_0)\grad p'$; linearized continuity:
  $\partial_t\rho'+\rho_0\divg\vel=0$; equation of state:
  $p'=c^2\rho'$.
  \item Taking the curl gives $\partial_t\curl\vel=0$, so transverse acoustic
  waves do not propagate in an ideal fluid.
  \item For the longitudinal part $\vel=\grad\phi$, the potential obeys
  $\lapl\phi=(1/c^2)\phi_{tt}$.
  \item The sound speed is $c^2=(\partial p/\partial\rho)_s$. For an ideal gas,
  $c=\sqrt{\gamma p_0/\rho_0}$; for a liquid,
  $c=\sqrt{B_s/\rho_0}$.
  \item Numerically, $c_{\mathrm{air}}\simeq343\,\mathrm{m/s}$ near room
  temperature, while $c_{\mathrm{water}}\simeq1.48\,\mathrm{km/s}$.
\end{itemize}

\subsection*{Further Reading}
\begin{itemize}
  \item L.\,D. Landau and E.\,M. Lifshitz, \textit{Fluid Mechanics}:
  sections on sound waves, compressibility, and viscosity-induced attenuation.
  \item D.\,J. Acheson, \textit{Elementary Fluid Dynamics}: discussion of small
  compressible disturbances and acoustic waves.
  \item G.\,K. Batchelor, \textit{An Introduction to Fluid Dynamics}: linearized
  compressible flow and sound-wave propagation.
\end{itemize}

\end{document}
